\documentclass[12pt, a4paper]{article}

% ---------- Packages ----------
\usepackage[margin=1in]{geometry}
\usepackage{graphicx}
\usepackage{amsmath}
\usepackage{booktabs}
\usepackage{longtable}
\usepackage{hyperref}
\usepackage{xcolor}
\usepackage{titlesec}
\usepackage{enumitem}
\usepackage{fancyhdr}
\usepackage{caption}
\usepackage{listings}
\usepackage{needspace}
\usepackage{times}

\hypersetup{
    colorlinks=true,
    linkcolor=blue,
    urlcolor=blue,
    citecolor=blue
}

\titleformat{\section}{\large\bfseries}{\thesection}{1em}{}
\titleformat{\subsection}{\normalsize\bfseries}{\thesubsection}{1em}{}

\setlength{\headheight}{15pt}
\pagestyle{fancy}
\fancyhf{}
\fancyhead[L]{Student Performance Prediction}
\fancyhead[R]{\thepage}
\renewcommand{\headrulewidth}{0.4pt}

\lstset{
    basicstyle=\ttfamily\footnotesize,
    breaklines=true,
    backgroundcolor=\color{gray!10},
    frame=single,
    columns=fullflexible
}

\title{}
\date{}

\begin{document}

% ================= TITLE PAGE =================
\begin{titlepage}
    \centering

    \vspace*{\fill}

    {\huge\bfseries\underline{STUDENT PERFORMANCE}}\\[0.35cm]
    {\huge\bfseries\underline{PREDICTION}}

    \vspace{2.5cm}

    {\large\bfseries Prepared by:}\\[0.4cm]
    {\large\bfseries Tanvir Kaur}\\[0.15cm]
    {\large (102303389)}

    \vspace*{\fill}

    \includegraphics[width=3.5cm]{ti.logo.png}\\[0.6cm]
    {\large\bfseries Thapar Institute of Engineering and Technology}\\[0.2cm]
    {\large\bfseries Patiala}

    \vspace{1.8cm}

    {\small\href{https://github.com/tanvirkaur30/student-performance-prediction}{github.com/tanvirkaur30/student-performance-prediction}}

    \vspace*{\fill}

\end{titlepage}

\tableofcontents
\newpage

% ================= ABSTRACT =================
\section*{Abstract}
Academic achievement is a cornerstone of a student's educational journey, influencing future
career opportunities and personal growth. Predicting student performance enables educational
institutions to identify and support students who may be at risk of underperforming. This
report presents a comprehensive machine learning pipeline — Logistic Regression, Decision
Tree, Random Forest, k-Nearest Neighbors, and XGBoost — applied to the UCI Student Performance
dataset. Beyond the baseline comparison, this iteration of the project adds engineered
features, stratified cross-validation, hyperparameter tuning, SHAP-based explainability, and
an ablation study that quantifies how much predictive power comes specifically from prior
academic grades versus demographic and behavioral factors. A tuned \textbf{XGBoost} model
achieved the best cross-validated performance (macro F1 $\approx 0.90$), and the full pipeline
is packaged as a reusable, deployable project with an interactive Streamlit demo.

\section{Introduction}
Student performance is a primary indicator of educational effectiveness and is critical for
forecasting future academic and professional success. Educational institutions increasingly
rely on data-driven approaches to identify students requiring intervention. Predictive
analytics powered by machine learning offers a robust framework to classify students by risk
level and design targeted support strategies.

Academic institutions often face challenges in predicting student success due to the
multifaceted nature of educational performance, influenced by cognitive, behavioral, social,
and psychological factors. Traditional analytical methods often fall short in capturing these
complexities. Machine learning, by contrast, offers a more sophisticated approach by detecting
hidden patterns and making data-driven predictions. This project explores the feasibility of
predicting student performance through machine learning methods, and — going beyond the
original coursework scope — investigates \emph{which} predictions are actually useful in
practice by separating the contribution of prior grades from other, earlier-available signals.

\subsection*{Scope of the Report}
This report investigates the application of supervised machine learning algorithms to predict
student performance using a dataset comprising demographic, academic, and behavioral
attributes. The goal is to provide actionable insights for educators, administrators, and
policymakers, enabling data-informed strategies to support student achievement, and to
demonstrate an end-to-end, production-style ML workflow (versioned pipeline, tuned models,
explainability, and a deployable interface).

\subsection*{Problem Statement}
Academic performance significantly affects students' life trajectories. However, predicting
performance is challenging due to the multitude of influencing factors such as demographics,
socio-economic background, behavioral habits, and psychological parameters. The core problem
addressed in this project is: \emph{``Can we predict students' academic performance accurately
based on available demographic, academic, and behavioral attributes — and how much of that
accuracy depends on already knowing their prior grades?''}

\subsection*{Objectives}
\begin{itemize}[noitemsep]
    \item Predict students' academic performance category (Low / Medium / High) from available attributes.
    \item Identify key determinants impacting academic success, with quantified feature importance.
    \item Evaluate and compare multiple ML models using robust cross-validation, not a single split.
    \item Tune the strongest candidates and quantify the gain from hyperparameter optimization.
    \item Separate the predictive contribution of prior grades from behavioral/demographic factors (ablation study).
    \item Provide actionable, explainable insights for educators via SHAP analysis.
    \item Package the work as a reusable, deployable, GitHub-ready project.
\end{itemize}

\section{Literature Review}
Recent years have seen a surge in research focused on predicting student performance using
machine learning and data mining techniques. Studies have explored a variety of algorithms,
including Decision Trees, Random Forests, Support Vector Machines, Naive Bayes, and Neural
Networks, each with varying degrees of success.

\textbf{Key findings from the literature:}
\begin{itemize}[noitemsep]
    \item \textbf{Feature Selection:} Prior academic performance, attendance, participation,
    and parental background are consistently identified as strong predictors.
    \item \textbf{Model Performance:} Ensemble methods like Random Forest and boosting
    algorithms often outperform single classifiers due to their ability to handle complex,
    non-linear relationships.
    \item \textbf{Explainability:} Recent work emphasizes model transparency, using techniques
    such as SHAP values to interpret feature importance and build stakeholder trust. This
    report directly applies SHAP to the final model (Section~\ref{sec:shap}).
    \item \textbf{Personalized Interventions:} Predictive models are increasingly used to
    design targeted interventions, such as personalized study recommendations and early
    warning systems.
\end{itemize}
Despite these advances, challenges remain in data quality, ethical considerations, and
generalizability across diverse educational contexts — points revisited in
Section~\ref{sec:limitations}.

\section{Methodology}

\subsection{Data Collection}
The project uses the UCI Student Performance Dataset (Mathematics course), comprising 395
students' demographic information, academic records, and lifestyle habits, collected from two
secondary schools in Portugal (Cortez \& Silva, 2008).

\textbf{Attributes included:}
\begin{itemize}[noitemsep]
    \item \textbf{Demographic:} age, sex, address, family size, parental education, job types.
    \item \textbf{Academic:} grades (G1, G2, G3), study time, failures, school support.
    \item \textbf{Behavioral:} absences, extracurricular activities, free time, alcohol consumption, health status.
\end{itemize}

\subsection{Feature Engineering}
Beyond the raw UCI columns, three engineered features were added to give models (and this
analysis) more structured signal:
\begin{itemize}[noitemsep]
    \item \texttt{avg\_alcohol} — mean of workday (\texttt{Dalc}) and weekend (\texttt{Walc}) alcohol consumption.
    \item \texttt{parent\_edu\_avg} — mean of mother's and father's education level.
    \item \texttt{avg\_prior\_grade} — mean of \texttt{G1} and \texttt{G2}, capturing overall prior trend.
\end{itemize}

\subsection{Data Preprocessing}
\begin{itemize}[noitemsep]
    \item \textbf{Missing values:} none present in the original dataset.
    \item \textbf{Encoding:} implemented via a single \texttt{scikit-learn} \texttt{ColumnTransformer}
    inside a \texttt{Pipeline} — one-hot encoding for binary and nominal categorical columns,
    standard scaling for numeric columns. This replaces ad-hoc \texttt{LabelEncoder} usage from
    the original version, which risked implying false ordinal relationships between nominal
    categories (e.g. parental job types) and could leak scaling statistics across CV folds.
    \item \textbf{Target construction:} \texttt{G3} (0--20) is bucketed into \emph{Low (0--9)},
    \emph{Medium (10--14)}, and \emph{High (15--20)}.
    \item \textbf{Splitting:} stratified 80/20 train-test split, with all model comparison done
    via stratified 5-fold cross-validation on the training set.
\end{itemize}

\subsection{Exploratory Data Analysis}
\begin{figure}[h]
    \centering
    \includegraphics[width=0.7\textwidth]{grade_distribution.png}
    \caption{Distribution of final grades (G3) across all 395 students.}
\end{figure}

\begin{figure}[h]
    \centering
    \includegraphics[width=0.6\textwidth]{class_balance.png}
    \caption{Class balance after bucketing G3 into Low / Medium / High. The dataset is
    moderately imbalanced (Low: 130, Medium: 192, High: 73), which is why macro-averaged
    F1-score — not raw accuracy — is used as the primary tuning and comparison metric.}
\end{figure}

Key EDA findings: strong positive correlation between \texttt{G1}, \texttt{G2}, and
\texttt{G3}; a negative relationship between absences and final grade; higher weekend alcohol
consumption associated with lower grades; and higher study time associated with better
outcomes.

\section{Model Building}

\subsection{Pipeline Design}
All five candidate models are wrapped in an identical \texttt{Pipeline(preprocessor, classifier)}
so that encoding and scaling are fit \emph{only} on training folds during cross-validation —
eliminating the data leakage risk present when preprocessing is done once on the full dataset
before splitting.

\subsection{Baseline Comparison (5-fold Stratified Cross-Validation)}
\begin{table}[h]
\centering
\begin{tabular}{lcc}
\toprule
\textbf{Model} & \textbf{CV Accuracy} & \textbf{CV Macro F1} \\
\midrule
Logistic Regression & 0.836 $\pm$ 0.019 & 0.834 $\pm$ 0.020 \\
Decision Tree        & 0.842 $\pm$ 0.022 & 0.844 $\pm$ 0.025 \\
Random Forest         & 0.867 $\pm$ 0.045 & 0.869 $\pm$ 0.042 \\
k-Nearest Neighbors   & 0.706 $\pm$ 0.050 & 0.696 $\pm$ 0.056 \\
\textbf{XGBoost}      & \textbf{0.886 $\pm$ 0.043} & \textbf{0.892 $\pm$ 0.038} \\
\bottomrule
\end{tabular}
\caption{Baseline 5-fold stratified CV results (untuned models, training set only).}
\end{table}

\begin{figure}[h]
    \centering
    \includegraphics[width=0.75\textwidth]{model_comparison.png}
    \caption{Cross-validated macro F1-score comparison across all baseline models.}
\end{figure}

XGBoost and Random Forest are the clear front-runners, confirming the literature's consistent
finding that ensemble/boosting methods outperform single classifiers on this type of
structured, moderately-sized tabular data.

\subsection{Hyperparameter Tuning}
The two strongest candidates were tuned using \texttt{GridSearchCV} (5-fold CV, scored on
macro F1):

\begin{itemize}[noitemsep]
    \item \textbf{Random Forest} — grid over \texttt{n\_estimators} $\in \{100,200,400\}$,
    \texttt{max\_depth} $\in \{\text{None},6,10,16\}$, \texttt{min\_samples\_split} $\in
    \{2,4,8\}$, \texttt{min\_samples\_leaf} $\in \{1,2,4\}$. Best: \texttt{n\_estimators=100},
    \texttt{max\_depth=None}, \texttt{min\_samples\_split=4}, \texttt{min\_samples\_leaf=1}
    $\to$ CV Macro F1 $= 0.891$.
    \item \textbf{XGBoost} — grid over \texttt{n\_estimators} $\in \{100,200,300\}$,
    \texttt{max\_depth} $\in \{3,4,6\}$, \texttt{learning\_rate} $\in \{0.05,0.1,0.2\}$. Best:
    \texttt{n\_estimators=300}, \texttt{max\_depth=4}, \texttt{learning\_rate=0.1} $\to$ CV
    Macro F1 $= 0.899$.
\end{itemize}

Tuned \textbf{XGBoost} was selected as the final model.

\section{Model Evaluation}

\subsection{Held-out Test Set Performance}
\begin{table}[h]
\centering
\begin{tabular}{lccc}
\toprule
\textbf{Class} & \textbf{Precision} & \textbf{Recall} & \textbf{F1-score} \\
\midrule
High   & 0.94 & 1.00 & 0.97 \\
Low    & 0.81 & 0.85 & 0.83 \\
Medium & 0.89 & 0.84 & 0.86 \\
\midrule
\textbf{Accuracy}   & \multicolumn{3}{c}{0.873} \\
\textbf{Macro avg}  & 0.88 & 0.90 & 0.89 \\
\textbf{Weighted avg} & 0.87 & 0.87 & 0.87 \\
\bottomrule
\end{tabular}
\caption{Classification report — tuned XGBoost on the held-out 20\% test set (79 students).}
\end{table}

\begin{figure}[h]
    \centering
    \includegraphics[width=0.6\textwidth]{confusion_matrix_best_model.png}
    \caption{Confusion matrix — tuned XGBoost, held-out test set. All 15 ``High'' students
    were correctly identified; most confusion occurs between the adjacent Low/Medium classes,
    which is expected given they are artificial buckets of a continuous grade.}
\end{figure}
\Needspace{8\baselineskip}
\subsection{Ablation Study: Predicting Without Prior Grades}
G1 and G2 dominate feature importance in essentially every student-performance study,
including this one — but that also makes the prediction less actionable: if a student's
period grades are already known, a model isn't required to notice they're struggling. The more
useful question is whether performance can be flagged \emph{before} G1/G2 exist, using only
demographic and behavioral data.

\begin{table}[h]
\centering
\begin{tabular}{lcc}
\toprule
\textbf{Feature set} & \textbf{Test Accuracy} & \textbf{Test Macro F1} \\
\midrule
With G1/G2 (full model)      & 0.899 & 0.915 \\
\textbf{Without G1/G2 (ablation)} & 0.557 & 0.473 \\
\bottomrule
\end{tabular}
\caption{Ablation study — Random Forest trained with vs.\ without prior-grade features.}
\end{table}

As expected, performance drops substantially when G1/G2 are
removed: test accuracy falls from 89.9\% to 55.7\%, a drop of
approximately 34 percentage points. However, the model still
performs above the majority-class baseline ($\approx 49\%$) using
only demographic, behavioral, and lifestyle features. This supports a two-stage early-warning design: a lighter model for
\emph{early-semester} risk flagging before grades exist, and the full model for
\emph{mid-year} intervention targeting once G1/G2 are available.

\section{Feature Importance and Explainability}
\label{sec:shap}
To move beyond raw feature-importance scores (which can be misleading for correlated or
one-hot-encoded features), SHAP (SHapley Additive exPlanations) values were computed for the
final XGBoost model.

\begin{figure}[h]
    \centering
    \includegraphics[width=0.8\textwidth]{shap_summary.png}
    \caption{SHAP summary plot — top 15 features by mean absolute SHAP value (averaged across
    the three output classes). Point color indicates the feature's value (red = high, blue =
    low) and horizontal position shows its impact on the model's output.}
\end{figure}

\textbf{Reading the plot:} \texttt{G2} and the engineered \texttt{avg\_prior\_grade} dominate,
confirming the correlation analysis. Beyond grades, weekend alcohol consumption
(\texttt{Walc}), absences, and age carry meaningful secondary signal — consistent with the
discussion in Section~\ref{sec:discussion} and the wider literature on behavioral predictors
of academic outcomes.

\section{Results Summary}
\begin{itemize}[noitemsep]
    \item \textbf{Best model:} Tuned XGBoost — CV Macro F1 $= 0.899$, Test Accuracy $= 87.3\%$, Test Macro F1 $= 0.888$.
    \item \textbf{Key predictors:} G2, avg\_prior\_grade, G1, absences, weekend alcohol use (Walc), age.
    \item \textbf{Ablation finding:} removing prior grades reduces
Random Forest test accuracy from 89.9\% to 55.7\%, a drop of
approximately 34 percentage points, demonstrating the strong
predictive contribution of prior academic performance.
    \item \textbf{Hyperparameter tuning gain:} CV Macro F1 improved from 0.892 (default XGBoost) to 0.899 (tuned) — a modest but consistent gain, showing the defaults were already reasonably close to optimal for this dataset size.
\end{itemize}

\section{Discussion}
\label{sec:discussion}
\begin{itemize}[noitemsep]
    \item \textbf{Academic history dominates:} Previous grades (G1, G2) are the strongest
    predictors, aligning with educational theory that past performance is a reliable indicator
    of future success.
    \item \textbf{Behavioral factors matter, but secondarily:} absences and alcohol use are
    significant contributors, but the ablation study shows they cannot substitute for the
    signal carried by prior grades.
    \item \textbf{Ensemble/boosting methods win:} XGBoost and Random Forest both outperform
    Logistic Regression, Decision Tree, and k-NN — confirming their suitability for this kind
    of structured, mixed categorical/numeric tabular data.
    \item \textbf{Practical implication:} a two-tier early-warning system is plausible — a
    lightweight model using only demographic/behavioral data early in the term, upgraded to
    the full model once period grades become available.
\end{itemize}

\section{Deployment}
To make this project usable beyond a notebook, the trained pipeline is wrapped in:
\begin{itemize}[noitemsep]
    \item \textbf{\texttt{src/predict.py}} — a reusable inference function and CLI for single-student predictions.
    \item \textbf{\texttt{app.py}} — an interactive Streamlit web app where a user adjusts sliders/dropdowns for a hypothetical student and receives a live Low/Medium/High prediction with class probabilities.
\end{itemize}
The full source, trained model, notebook, and this report are published at:
\url{https://github.com/tanvirkaur30/student-performance-prediction}

\section{Limitations}
\label{sec:limitations}
\begin{itemize}[noitemsep]
    \item The dataset is limited to 395 students from two schools in Portugal, which may limit generalizability to other educational systems.
    \item Psychological and socio-emotional factors were not captured in the original data and could further improve prediction accuracy.
    \item The Low/Medium/High bucketing is a modeling choice; a regression formulation of G3 was not explored here but is a natural extension.
    \item Ethical considerations — data privacy, and the risk of a predictive label influencing how a student is treated by staff — must be addressed before any real-world deployment.
\end{itemize}

\section{Conclusion}
Predicting student performance using machine learning is highly feasible: a tuned XGBoost
model reached 87.3\% test accuracy and 0.888 macro F1 on a held-out set. Prior academic
grades remain the dominant predictor, but the ablation study shows meaningful signal persists
in purely demographic and behavioral data, supporting early-intervention use cases. SHAP
analysis confirms the model's reasoning aligns with domain intuition (grades, attendance, and
lifestyle factors), which is important for building institutional trust in such a system.

Future work could explore regression on the raw G3 score, incorporate mental-health or
online-learning-behavior indicators, and validate the approach across additional schools or
countries before any operational deployment.

\section*{References}
\begin{itemize}[noitemsep]
    \item Cortez, P., \& Silva, A. (2008). \emph{Using Data Mining to Predict Secondary School Student Performance}. University of Minho, Portugal.
    \item UCI Machine Learning Repository. \emph{Student Performance Data Set}. \url{https://archive.ics.uci.edu/ml/datasets/Student+Performance}
    \item James, G., Witten, D., Hastie, T., \& Tibshirani, R. (2013). \emph{An Introduction to Statistical Learning}. Springer.
    \item Pedregosa, F. et al. (2011). Scikit-learn: Machine Learning in Python. \emph{Journal of Machine Learning Research}, 12, 2825--2830.
    \item Chen, T., \& Guestrin, C. (2016). XGBoost: A Scalable Tree Boosting System. \emph{Proceedings of the 22nd ACM SIGKDD}.
    \item Lundberg, S. M., \& Lee, S.-I. (2017). A Unified Approach to Interpreting Model Predictions. \emph{Advances in Neural Information Processing Systems}.
    \item Han, J., Kamber, M., \& Pei, J. (2011). \emph{Data Mining: Concepts and Techniques}. Elsevier.
    \item Kuhn, M., \& Johnson, K. (2013). \emph{Applied Predictive Modeling}. Springer.
\end{itemize}

\end{document}